Segment-onset delay is a natural timing signal in simultaneous interpreting, but it is not equivalent to a professional evaluator's judgment that an interpretation follows the source promptly. We formulate professional promptness-rating prediction as a supervised segment-level task using source text, interpreted output, and timing information. Our dataset contains \CohortN\ jointly rated Chinese--English and English--Chinese segments from \SpeechGroupN\ source-speech groups and \InterpreterN\ interpreters, evaluated for language quality (LQ), expressiveness (EXP), and promptness by two senior professional interpreters. We use their predefined two-rater professional aggregate, the equal-weight mean, while auditing evaluator-specific calibration. Under cross-fitted source-speech-group-held-out evaluation, segment-onset delay alone achieves Pearson $r=\DelayPearson$; predicted LQ and EXP combined with delay reach $r=\QualityDelayPearson\pm\QualityDelayPearsonSD$. A stronger deterministic lexical/structural-plus-delay baseline reaches $r=\StructureDelayPearson$, and the full model reaches $r=\FullPearson\pm\FullPearsonSD$, indicating a modest, heterogeneous complementary association from predicted quality beyond structure and timing. Human-label cross-rater analyses show positive but asymmetric transfer, indicating that the quality--promptness relationship is not confined to within-evaluator score coupling but remains evaluator-dependent. Interpreter-disjoint evaluation remains positive but heterogeneous. At camera-ready, we will release anonymized labels, delay records, rubric, fold manifests, and code subject to the study's governance constraints.
